\section{Cronograma de Actividades}
\label{sec:cronograma}

El proyecto está diseñado para desarrollarse en un periodo de 8 semestres (4 años), distribuidos de acuerdo con las etapas clave de investigación, análisis y difusión de resultados descritas en la Tabla~\ref{tab:cronograma}.

\begin{table}[htbp]
    \centering
    \caption{Cronograma semestral de actividades del proyecto de investigación doctoral.}
    \label{tab:cronograma}
    \resizebox{\textwidth}{!}{%
    \begin{tabular}{lcccccccc}
        \toprule
        \textbf{Actividad / Tarea} & \textbf{S1} & \textbf{S2} & \textbf{S3} & \textbf{S4} & \textbf{S5} & \textbf{S6} & \textbf{S7} & \textbf{S8} \\
        \midrule
        Revisión bibliográfica y actualización teórica & X & X & X & X & X & X & X & X \\
        Definición y modelado de interfaces Heusler/MXene & X & X & & & & & & \\
        Optimización geométrica y convergencia (DFT) & & X & X & & & & & \\
        Cálculos de estructura de bandas, DOS y magnetismo & & & X & X & & & & \\
        Cálculos relativistas (SOC) y análisis topológico & & & & X & X & X & & \\
        Simulación de modulación por strain y campos externos & & & & & X & X & & \\
        Escritura de manuscritos científicos (artículos) & & & & X & X & X & X & \\
        Redacción del documento de tesis y revisión & & & & & & & X & X \\
        Preparación y defensa del examen de grado & & & & & & & & X \\
        \bottomrule
    \end{tabular}%
    }
\end{table}

\noindent \textbf{Detalles de los hitos principales:}
\begin{itemize}
    \item \textbf{Fin del Año 1 (S2):} Modelos de interfaz listos y optimizados estructuralmente.
    \item \textbf{Fin del Año 2 (S4):} Propiedades magnéticas y electrónicas calculadas (envío de primer artículo).
    \item \textbf{Fin del Año 3 (S6):} Fases topológicas caracterizadas y mapeo del efecto de perturbaciones externas completado (envío de segundo artículo).
    \item \textbf{Fin del Año 4 (S8):} Redacción completa del borrador de tesis y examen de defensa doctoral.
\end{itemize}
